% The evaluation loop that isolates the advantage of prediction (P1 analog: fig_architectures).
% Purely structural (no counts): shows where a benchmark must run BOTH a direct-VLA policy and a
% world-model policy under one closed loop, and where the per-capability advantage is measured.
\begin{figure}[t]
\centering
\resizebox{\textwidth}{!}{%
\begin{tikzpicture}[
 box/.style={draw=black!55, rounded corners, align=center, font=\footnotesize\bfseries,
             minimum height=9mm, text width=20mm, inner sep=3pt},
 io/.style={draw=none, font=\footnotesize\bfseries\itshape, align=center},
 gap/.style={box, fill=secEl, draw=secE!60!black, line width=1.1pt, text width=24mm},
 ar/.style={-{Stealth[length=2mm]}, darkgray, line width=0.8pt},
 fb/.style={-{Stealth[length=1.8mm]}, gray!70, line width=0.6pt},
]
 \node[io] (task) at (0,0) {task /\\ scene};
 \node[box, fill=cat-vla] (vla) at (3.1,1.15) {direct VLA\\ policy};
 \node[box, fill=secFl] (wm) at (3.1,-1.15) {world model:\\ predict $\rightarrow$ plan};
 \node[box, fill=secCl] (env) at (6.4,0) {environment\\ (closed loop)};
 \node[box, fill=secBl] (succ) at (9.4,0) {task\\ success};
 \node[gap] (adv) at (12.5,0) {advantage of\\ prediction ($\delta$)\\ per capability ($\beta$)};

 \draw[ar] (task) -- (vla);
 \draw[ar] (task) -- (wm);
 \draw[ar] (vla) -- (env);
 \draw[ar] (wm) -- (env);
 \draw[ar] (env) -- (succ);
 \draw[ar] (succ) -- (adv);
 % state-feedback loops (what makes it closed-loop); label once on top
 \draw[fb] (env.north) to[out=90,in=90] node[above,font=\scriptsize\bfseries,text=gray!20!black]{state feedback} (vla.north);
 \draw[fb] (env.south) to[out=-90,in=-90] (wm.south);
 % gap callout: the two ways today's suites break the loop
 \node[draw=BrickRed!55!black, dashed, rounded corners, fill=BrickRed!5, inner sep=4pt,
       text width=30mm, font=\scriptsize\bfseries, text=BrickRed!85!black, align=center] at (3.1,-3.1)
   {model-agnostic suites run \textbf{one} branch (no contrast); world-model suites never reach the loop};
 \draw[BrickRed!55!black, dashed, -{Stealth[length=1.6mm]}] (3.1,-2.2) -- (3.1,-2.55);
\end{tikzpicture}}
\caption{\textbf{The evaluation loop that isolates the advantage of prediction.} A benchmark that
can answer ``does world modelling help a robot act'' must run a direct Vision-Language-Action policy
\emph{and} a world-model policy (predict then plan) under one closed loop, then report the
closed-loop gain ($\delta$) sliced by capability ($\beta$). Today's suites break this loop in one of
two ways: model-agnostic policy suites host a single policy and never build the contrast, while
world-model suites score the prediction open-loop and never execute it (\S\ref{sec:gap}).}
\Description{A left-to-right diagram. A task/scene forks into two branches, a direct VLA policy and a
world model that predicts then plans; both feed an environment closed loop with state feedback,
producing task success, which feeds an advantage-of-prediction measurement sliced per capability. An
annotation notes that model-agnostic suites run only one branch and world-model suites never close
the loop, so the contrast is missing.}
\label{fig:eval_loop}
\end{figure}
